% ==============================
% Vorspann des LaTeX Dokumentes
% ==============================

% --- Definition der Dokumentenklasse: \documentclass[Optionen]{Klassenname}
% --- Dokumentenklasse gibt zahlreiche Vereinbarungen ueber Layout und Gliederung vor
% --- Wir verwenden hier die KOMA-Script Klasse 'srcreprt'
% --- Siehe dazu Anleitung zum KOMA-Script, insb. Abschnitt 2.6, Kapitel 3,
\documentclass[%
	a4paper,%							Papiergroesse im DIN A4-Format
	oneside,%							Einseitiger Druck
%	twoside,%							Zweiseitiger Druck
	12pt,%								Schriftgroesse
%	draft,%								Option unterscheidet zw. Dokument im Entwurfsstadium und fertigen Dokument
	toc=listof,%					Abbildungs- und Tabellenverzeichnis erhalten einen Eintrag im Inhaltsverzeichnis ohne Nummerierung
	toc=bibliography,%		Literaturverzeichnis erhaelt einen Eintrag im Inhaltsverzeichnis ohne Nummerierung
%	toc=flat,%						Inhaltsverzeichnis erhaelt eine tabellarische Form (ohne Einruecken der unteren Ebenen)
%	chapterprefix,%				Kapitelueberschriften werden nicht in kompakter Form dargestellt
%	appendixprefix%				Anhangueberschriften werden nicht in kompakter Form dargestellt
%	titlepage%						Titel und Zusammenfassung erscheinen auf einer eigenen Seite (Voreinstellung f�r 'book' und 'report')
%	fleqn%								Mathematische Formeln werden linksbuendig statt zentriert gesetzt
%	leqno%								Mathematische Formeln werden auf der linken Seite statt auf der rechten Seite nummeriert
]{scrreprt}%


% --- Eigene Stildatei (Hier werden alle grundlegenden Einstellungen fuer Layout und Packages vorgenommen)
% --- Moeglicher Inhalt: Einbindung von haeufig verwendeten Pakete und Erweiterungspaketen
% ---										 Definition eigener Befehle
% ---                    Ausfuehrung global wirkender Befehle (z.B. globale Layout Einstellungen etc.)
\usepackage{styles/abschlussarbeit}


% =============================
%	Hauptdokument mit Textinhalt
% =============================

% --- Aufteilung in mehrere Quelldateien
%			LaTeX unterstuetzt die Aufsplittung in Teildokumente, so dass in einem Uebersetzungslauf nur ein gewisser Anteil verarbeitet wird
%			\include{Dateiname} - bindet 'Dateiname' (ohne .tex Endung) in das Dokument, LaTeX sorgt fuer Seitenumbruch vor und nach dem Inhalt dieser Datei und verwaltet eine eigene Hilfsdatei ('Dateiname'.aux) mit Querreferenzen u.ae.
%			\includeonly{Dateiname,Dateiname,...} - werden nur diejenigen Dateien tatsaechlich mit \include eingebunden, die hier aufgaehlt werden. Ist \includeonly nicht vorhanden, so werden alle Dateien eingebunden.

%\includeonly{chapters/heranfuehrende_kapitel,chapters/hauptteil}

\begin{document}

% --- Titelblatt
% ===========
% Titelseite
% ===========

% --- Kopf- und Fusszeilen der Titelseite sind leer
\thispagestyle{empty}

%\hspace*{1em}
\vfill
	\begin{center}
			\includegraphics[width=10cm]{images/TU_Logo}
		\par
		\vspace*{8ex}
\Huge
%			Titel der Abschlussarbeit
			\textsc{Model order reduction}
		\par
\normalsize
		\vspace*{8ex}
\large
			\textsc{lecture notes}\\
		\par
\normalsize
		\vspace*{6ex}
			Technische Universit�t Berlin\\
			Fakult�t II: Mathematik\\
			Fortgeschrittene Themen der Numerik
	\end{center}
	\par
	\vspace*{12ex}
		Dozent:\\
\large
			\textit{Dr. Schr�der}
	\par
\normalsize
	\vspace*{2ex}
		Written by:\\
\large
			\textit{Ines Ahrens}
	\par
\normalsize
	\vspace*{4ex}
		Berlin, Oktober 2015
\vfill
\hspace*{1em}

% kleine roemische Seitenzahlen im Vorspann der Arbeit, Seitenzaehler wird automatisch auf 1 gesetzt

% --- Inhaltsverzeichnis wird automatisch erzeugt
\tableofcontents

% technische Gruende, damit Inhaltsverzeichnis roemische Seitenzahl erhaelt
\newpage
% arabische Seitenzahlen im Rest der Arbeit, Seitenzaehler wird automatisch auf 1 gesetzt
\pagenumbering{arabic}

% --- Einleitung (Introduction)
% ===========
% Einleitung
% ===========

\chapter{Introduction: What is model order reduction}
\label{chap:introduction}

\begin{itemize}
	\item rigorous definition is difficult/ impossible
	\item MOR is the approximation of large-scale systems of equations by small scale ones with similar solution properties
\end{itemize}

Here, we consider the MOR of dynamical systems of the form
\begin{align*}
	& \dot{x}(t) : = f(x(t), u(t), t) \hspace{2ex} x(t_0)=x_0 \hspace{2ex} t \in I \subset \R \hspace{2ex} I=[t_0, t_f] \\
	& y(t)=g(x(t),u(t),t) 
\end{align*}
with 
\begin{align*}
	& x(\dot ) : I \rightarrow \R^n & \text{state (function)}\\
	& u(\dot ) : I \rightarrow \R^m & \text{input (function)}\\
	& y(\dot ) : I \rightarrow \R^p & \text{output (function)}\\
	& f : \R^n \times \R^m \times I  \rightarrow \R^n & \text{dynamics of system}\\
	& g: \R^n \times \R^m \times I \rightarrow \R^p & \text{output map}\\
\end{align*}

\begin{note}
	if u is fixed we can calculate u, then y. In this lecture we only consider systems where the pde is solvable. 
\end{note}

usually the input can be chosen by a system operator in order to influence the system state. The output consists of the system parameters that are servable or that are of interest to the operator.

The task of MOR can be described as 
\begin{itemize}
	\item find functions 
	\begin{align*}
		& \tilde{f} : \R^l \times \R^m \times I  \rightarrow  \R^n \\
		& \tilde{g}: \R^l \times \R^m \times I \rightarrow \R^p
	\end{align*}
	with $l <<n$ and $\tilde{x}_0 \in \R^l$ such that the output $\tilde{y}$ of the small  system 
	\begin{align*}
		& \dot{\tilde{x}}(t) : = f(\tilde{x}(t), u(t), t) \hspace{2ex} \tilde{x}(t_0)=x_0 \\
		& \tilde{y}(t)=g(\tilde{x}(t),u(t),t) 
	\end{align*}
	is very close to $y$ for most (or all admissable) input functions $u:I \rightarrow \R^m$
	\item meaning of very close will become clear later
	\item purpose of MOR is to aid numerical simulation of large scale systems, which is computationally very demanding for large $n$, but cheap for small $l$. 
\end{itemize}

\begin{exa}
	Consider the linear system with $I=(0,\infty]$
	\begin{align*}
		\begin{pmatrix}
			\dot{x}_1(t) \\\dot{x}_2(t) 
		\end{pmatrix}
		= 
		\begin{pmatrix}
			-1 & 0 \\ 0 & -2
		\end{pmatrix}
		\begin{pmatrix}
			x_1(t) \\ x_2(t)
		\end{pmatrix}
		+
		\begin{pmatrix}
			1 \\ \frac{1}{10}
		\end{pmatrix}
		u(t), \hspace{2ex} 
		\begin{pmatrix}
			x_1(0) \\ x_2(0)
		\end{pmatrix}
		= \begin{pmatrix}
			0 \\ 0
		\end{pmatrix}
	\end{align*}
	
	The state variable $x_1(t)$, $x_2(t)$ are mostly decoupled, so using the formular of variation of the constant, we have
	\begin{eqnarray*}
		y(t)=2x_1(t)+x_2(t) & = 2 \int\limits_0^t e^{- (t- \tau)} u(\tau) \diff \tau + \frac{1}{10} \int\limits_0^t e^{-2(t - \tau )} u(\tau) \diff \tau \\
		& \sim 2 \int\limits_0^t e^{-(t-\tau)} u(\tau) \diff \tau 
	\end{eqnarray*}
	The term after $1/10$ can be erased. 
	
	Hence, the smaller System  $\dot{\tilde{x}}(t)= - \tilde{x}(t) +u(t), \hspace{2ex} \tilde{x}(0)=0, \tilde{y}(t)=2 \tilde{x}(t)$ is an approximation of the original one. 
\end{exa}

\begin{exa}
	We consider the temper distribution $T(t, \xi)$ in a beam of length 1. (\dots) in the lecture we calculate this discrete and the solution $u$ is a \glqq Schock \grqq
\end{exa}

\begin{note}
	lecture 2 is missing
\end{note}

\section{?}

\section{?}

\textbf{recap: dynamical Systems }

\begin{align*}
		& x(\dot ) : I \rightarrow \R^n & \text{state (function)}\\
		& u(\dot ) : I \rightarrow \R^m & \text{input (function)}\\
		& y(\dot ) : I \rightarrow \R^p & \text{output (function)}
\end{align*}

LTI system: There are matrices $A,B,C,D$ s.t. 
\begin{align}
	\label{eq:LTI}
	& \dot{x}(t)=A x(t) + B u(t), \hspace{2ex} x(0)=x_0\\
	& y(t)=C \dot{x}(t) + D u(t)
\end{align}


An LTI can be given by
\begin{itemize}
	\item  $A,B,C,D, x_0$. Then $y(t)= C e^{A(t-\tau)} x_0 + \int\limits_0^t C e^{A(t-\tau)}  B u(\tau) \diff \tau + D u(t)$\item impulse response $y_d$. then $y=y_d * u_0$ (only for $x=0$)
	\item transfer function G. then $L(y)(s)=\hat{y}(u)= G(s)\hat{u}(s)=L(u(s))$ (only for $x=0$)
\end{itemize}

the three parametrization can be transformed into one another. 
\begin{align*}
	y_0(t) & = \left\{ \begin{array}{ll}
			C e^{At}B+ \partial (t) D & t \leq 0 \\
			0	& t<0
	\end{array} \right. \\
	G(s) & = C(sT-A)^{-1} B+D \\
	G(s) & = L(y_d)(s)
\end{align*}
For $u(t)=e^{iwt}I_m$ and $x_0=(iwI-A)^{-1}B$ we have $y(t)= \cdots = G(iw)u(t)$ 
So, excitation of a suitably initialized LTI-System \ref{eq:LTI} with an input of frequency $\omega$ leads to an output of frequency $\omega$ only amplified by a (frequency dependent) factor $G( i \omega) $. For that reason the distraction of G to the imaginary axis is called frequency response.

\section{realization} 

The entries of the matrix valued function G must be rational functions. To see this we use

\begin{defi}
	For $A \in \C^{n \times n}$ we define its adjoint $adj(A) \in \C^{n \times n}$ as
	\begin{align*}
		(adj(A))_{ij}= (-1)^{ij} \det A^{ij}
	\end{align*}
	where $A^{ij}$ is the matrix so that the i-th column and the j-th row is removed. 
\end{defi}

\begin{thm}
	For invertible $A \in \C^{n \times n } $ we have 
	\begin{align*}
		A^{-1}= \frac{adj(A)}{\det A}
	\end{align*} 
\end{thm}
\begin{proof}
	apply Cramers rule to $Ax=I$
\end{proof}

Hence, for the transfer function we have 
\begin{align}
\label{eq:N(s)}
	G(s)=C(sI-A)^{-1}B+D= C \frac{adj(sI-A)}{\det (sI-A)} B + D \notag \\
	= \frac{adj(sT-a)B+\det (sI-A) D}{\det (sT-A)}= \frac{N(s)}{d(s)}
\end{align}
where $d(s)=\det(sI-A)$ is a polynomial of degree n and $N(s):= adj(sI-A)B+\det(sT-A)D$ is a $p \limes m$ matrix polynomial. 

Matrix polyonmials are a matrix whose entries are polynomials. 

\begin{defi}
	The set of $p \times m$ matrix-valued rational functions G st $G(s) \in \R^{p \times m}$ for almost all $s \in \R $ is denoted by $\R^{p \times m}(S)$. Moreover $G \in R^{p \times m}(S)$ is called (strictly) proper if \begin{align*}
		\lim\limits_{s \rightarrow \infty} G(s)
	\end{align*}
	exists (or is zero) 
\end{defi}

\begin{align*}
	\R^{p \times m}_p(s):= \{G \in \R^{p \times m}(s)| G \text{ is proper} \} \\
	\R^{p \times m}_{ps}(s):= \{G \in \R^{p \times m}(s)| G \text{ is strict proper} \} 
\end{align*}

Since 
\begin{align*}
	\lim\limits_{s \rightarrow \infty} (sI -A)^{-1} =0
\end{align*}
we have $G(s) \rightarrow D$.

So a transfer function is proper. It's even strictly proper if $D=0$. Thus the degree of the numerator $N(s)$ in \ref{eq:N(s)} is at most n. Also \ref{eq:N(s)} every pole of $G(s)$ must be a zero of a i.e. an eigenvalue of $A$. The converse is not true. 

\begin{exa}
	\begin{align*}
		A=\begin{pmatrix}
		1 & 1 \\ 0  & 1 
		\end{pmatrix}
		, \hspace{2ex}
		B = \begin{pmatrix}
		3 \\ 0
		\end{pmatrix}, \hspace{2ex}
		C=(1,1), \hspace{2ex} D=0 
	\end{align*}
	So,
	\begin{align*}
		G(s)=\frac{3}{s-1}
	\end{align*}
	The sole pole ist at one which is an eigenvalue, but the eigenvalue 2 is not a pole. Also for $A=1, B=3, C=1, D=0$ we have $G(s)=\frac{3}{s-1}$
\end{exa} 

\begin{defi}
	For given $G \in \R^{p \times m}_p(s)$ the quadruple $[A,B,C,D] \in \R^{n \times n} \times \R^{n \times n} \times \R^{p \times n} \times \R^{p \times m}$ is called a realization (of G), if $G(s)=C(sI-A)^{-1}B + D $. A realization is called minimal if n is minimal among all realizations. Two realizations $[A_1,B_1,C-1,D_1]$ and $[A_2,B_2,C_2,D_2]$  are called equivalent if there is an invertible $T \in \R^{n \times n}$ s.t. $[A_2,B_2,C_2,D_2]= [T^{-1}A_1 T, T^{-1}B_1,C_1T,D_1]$
\end{defi}

\begin{thm}
	Equivalent realizations have the same transfer function. 
\end{thm}

\begin{thm}
	Every $G \in \R^{p \times m}_p(s)$ have a realization. 
\end{thm}



% --- Hauptteil der Arbeit I
% =======================
% Hauptteil der Arbeit I
% =======================

\chapter{Gruppenentscheidungen}

\section{Individualpräferenz und Gruppenentscheidungen}

Sei A eine endliche Menge von Kandidaten.

\begin{defi}[Präferenz]
	Sei A eine endliche Menge von Kandidaten. Einr Präferenz entsteht, wenn jeder Wähler jeden Kandidaten $x \in A$ eine natürliche Zahl als Rang $r(x)$ zuordnet. Bei zwei Kandidaten $x, y \in A$ bedeutet $r(x)<r(y)$, dass $X$ gegenüber $y$ bevorzugt wird. 
\end{defi}

\begin{defi}[Rangabbildung]
	Eine Rangabbildung ist eine surjektive Abbildung $r$ von der Kandidaten Menge A auf einen Abschnitt $\{ 1, \dots, k \} \subset \N $ 
	\begin{align*}
		r:A \rightarrow \{ 1, \dots, k \} 
	\end{align*}
\end{defi}

\begin{defi}[Relation]
	Eine Relation $R$ auf der Menge A heißt $\forall x,y,z \in A$
	\begin{itemize}
		\item transitiv, falls $x Ry, yRz \Rightarrow xRz $
		\item symmetrisch, falls $xRy \Rightarrow yRx$
		\item reflexiv $xRx$
		\item Quasiordnung $R$ ist reflexiv und transitiv
		\item asymetrisch $xRy \Rightarrow \neq (yRx)$
		\item konvex: jede zwei Elemente sind vergleichbar:$xRy \vee yRx$
	\end{itemize}
\end{defi}

\begin{defi}[Präferenzrelation]
	Die Rangabbildung r definiert eine Präferenzrelation $\rho \subset A \times A$ über 
	\begin{align}
		\label{eq:rangabbildung}
		x \rho y : \Leftrightarrow r(x) < r(y)
	\end{align}
\end{defi}
$\rho$ ist transitiv und asymetrisch. 
\begin{defi}
	Sei $\rho$ die durch der Rangabbildung $r$ definierte Präferenzrelation.  
	Die Menge aller so darstellbaren Relationen nennen wir 
	\begin{align*}
		\mathcal{P}_A:=\{ \rho \subset A \times A| \rho \text{ erfüllt \eqref{eq:rangabbildung} für eine Rangabbildung } r  \}
	\end{align*}
\end{defi}

Wir können auch noch Paare hinzunehmen, die beide den gleichen Rang haben: 
\begin{align}
	\label{eq:rangabb_stern}
	x \rho* y : \Leftrightarrow r(x) \leq r(y)
\end{align}
$\rho*$ ist transitiv, reflexiv und konvex. Es gilt $\rho \subset \rho* \subset A \times A$
und damit definieren wir 
\begin{align*}
	\mathcal{P}_A^*:=\{ \rho^* \subset A \times A| \rho^* \text{ erfüllt \eqref{eq:rangabb_stern} für eine Rangabbildung } r \}
	\end{align*}

\begin{thm}
	Sei $\rho$ die Präferenzrelation und $\rho*$ zugehörig. Dann gilt: 
	\begin{align*}
		x \rho y \Leftrightarrow \neq (x \rho* y)
	\end{align*}
\end{thm}
\begin{noti}
	in Vorlesung ohne Beweis. Vll selber machen
\end{noti}
Mit diesem Theorem kann man auch ohne die Rangabbildung $r$ zu kennen, aus gegebenen $\rho$ das zugehörige $\rho*$ bestimmen.  
 
Die Relationen sind also zueinander invers komplementär. 

\begin{defi}[Kollektive Auswahlfunktion]
	Sei $I=\{1, \dots, n\}$ eine Wählermenge und $\mathcal{P}_A$ wie oben definiert. Die kollektive Auswahlfunktion ist gegeben durch
	\begin{align*}
		K: \mathcal{P}_A^n \rightarrow \mathcal{P}_A
	\end{align*}
\end{defi}
Dabei repräsentiert jedes $\mathcal{P}_A$ aus der Eingabemenge eine Präferenz eines Wählers. Die Kollektive Auswahlfunktion berechnet die Präferenz der Gesamtheit. 

Wir haben zwei wesentliche Bedingungen an die Entscheidungsfunktion
\begin{itemize}
	\item Die Abbildung $K$ muss total sein, d.h. jedem Element des Definitionsbereiches muss ein Bild zugeordnet werden.
	\item Das Ergebnis der Wahl muss ebenfalls in $\mathcal{P}_A$ sein. 
\end{itemize} 

\section{Beispiele für Entscheidungsverfahren}
Das Ziel dieses Kapitels ist es, eine kollektive auswahlfunktion zu konstruieren, die besonders gerecht ist. 
Sei in diesem Kapitel $I=\{1, \dots, n\}$ die Menge von Wählern mit den Präferenzrelationen $\rho_i$ $\forall i \in I$. 

\subsection{externer Diktator}
Es gilt für beliebiges aber festes $\rho_E$: 
\begin{align*}
	K_{\rho_E}^E(\rho_1, \dots, \rho_n)= \rho_E
\end{align*}
Die ist offfensichtlich eine Abbildung

\subsection{interner Diktator}
Es gibt beim externen Diktator einen Wähler $d \in I$ mit 
\begin{align*}
	K_d^D(\rho_1, \dots, \rho_n)= \rho_d
\end{align*}

\subsection{Rangaddition}
Zu jeder Individualpräferenz $\rho_i$ gibt es eine Rangabbildung $r$. Bei der Rangaddition werden die Kandidaten 
nach der Summe ihrer Rangzahlen bewertet. Es ist $K^A(\rho_1, \dots, \rho_n)$ die Relation $\rho$ mit
\begin{align*}
	x \rho y : \Leftrightarrow \sum_{i=1}^n r_i(x) < \sum_{i=1}^n r_i(y)
\end{align*}
Das Problem ist, dass dies i.A. keine Rangabbildung ist, da die Abbildung nicht surjektiv sein muss (z.B. kein Kandidat hat genau 5 Punkte). Eine mögliche Lösung ist es, die nicht erreichten Punkte zu streichen und die anderen Kandidaten aufrücken zu lassen. 

\subsection{Cordocct Verfahren} 
Dieses Verfahren vergleicht die Kandidaten paarweise miteinander. Die kollektive Präferenzrelation lautet damit:
\begin{align*}
	x \rho y :\Leftrightarrow |\{i\in I | x \rho_i y \} | > |\{i \in I | y \rho_i x\} |
\end{align*} 

\subsection{Einstimmigkeit}
Setze Kandidat $x$ vor $y$, falls alle Wähler $x$ vor $y$ setzen. Die Kollektive Auswahlpräferenz lautet:
\begin{align*}
	x \rho y : \Leftrightarrow \forall i \in \{1, \cdots n\} x \rho_i y
\end{align*}
Falls nur ein Wähler $x$ genau so schätzt wie $y$,gilt $y \rho^* x$  .
Deshalb ist dieses Verfahren in der Praxis kaum anwendbar. Außerdem muss das Ergebnis für die Kandidatenmenge $A$ mit $|A|>2$ nicht unbedingt in $\mathcal{P}_A$ sein, da nicht unbedingt jeder Kandidat plaziert werden kann. 

\subsection{Borda-Wahl}  
Jeder Wähler gibt den k Kandidaten Punkten. Der erst präferierte von Wähler $i$ erhält $k$ Punkte, der letzte 1 Pkt. Der Kandidat mit den meisten Punkten gewinnt. 

\section{Bedingungen an die Auswahlfkt und Satz von Arrow}

Wir wollen zwei Bedingungen an ein gerechtes Verfahren stellen. 

\textbf{Parteo Bedingung: } 
Die Gesamtheit für einen beliebigen Kandidaten kann durch eine Einstimmigkeit jede gewünschte Richtung erreichen.

\textbf{Unabhängigkeit von irrelevanten Ereignisssen: } Die Reihung zwischen zwei Kandidaten kann nicht dadurch geändert werdne, dass Wähler ihre Präferenz bzgl. eines dritten Kandidaten verändern. 

\begin{defi}[Pareto Bedingung]
	Eine kollektive Auswahlfunktion $K: \mathcal{P}_A^n \rightarrow \mathcal{P}_A$ erfüllt die Pareto Bedingung, falls für alle $\rho_i \in \mathcal{P}_A, i \in I$ mit $\rho = K(\rho_1, \dots, \rho_n)$ und für alle Kandidaten $x,y \in A$ gilt: 
	\begin{align*}
		(\forall i \in O: x \rho_i y) \Rightarrow x \rho y
	\end{align*}
\end{defi} 

Die Pareto Bedingung kann auch für $\rho^*$ definiert werden, nur ein Verfahren, das jeden Kandidaten immer gleich bewertet, würde die Bedingung auch erfüllen. 

\begin{defi}[Unabhängigkeit von irrelevanten Alternativen/Ereignissen]
	Eine kollektive Auswahlfkt $K: \mathcal{P}_A^n \rightarrow \mathcal{P}_A$ erfüllt die unabhängigkeit von irrelevanten Ereignissen, wenn 
	\begin{align*}
		\forall \rho_i, \rho_i' \in \mathcal{P}_A, i \in I
	\end{align*}
	mit $\rho=K(\rho_1, \dots, \rho_n)$ und $\rho'=K(\rho_1', \dots, \rho_n')$ und für alle $x,y \in A$ gilt: 
	\begin{align*}
		(\forall i \in I \text{ mit } x \rho_i y \Leftrightarrow x \rho_i' y) \Rightarrow (x \rho y \Rightarrow x \rho' y)
	\end{align*}
\end{defi}

In Worten bedeutet das: Wenn kein Wähler seine Meinung bzgl x und y ändert, dann solte auch im Ergebnis die Reihung von x und y gleich bleiben. 

Nur noch der interne Diktator erfüllt alle Bedingungen einschließlich der Pareto Bedingung und der unabhängigkeit von irrelevanten Ereignissen. 

\begin{thm}[Satz von Arrow]
	Es sei A mit $|A|>2$ eine Kandidatenmenge und $K:  \mathcal{P}_A^n \rightarrow \mathcal{P}_A$ eine kollektive Auswahlfkt, die die Pareto Bedingung erfüllt und die Unabhängigkeit von irrelevanten Ereignissen erfüllt. Dann gilt es einen internen Diktator. D.h.
	\begin{align*}
		\exists d \in I: \forall (\rho_1, \dots, \rho_n) \in \mathcal{P}_A^n : \forall (x,y) \in A \times A \hspace{1ex} x \rho_d y \Rightarrow x \rho y
	\end{align*}
	mit $\rho = K(\rho_1, \dots, \rho_n)$
\end{thm}
\begin{proof}
	Seehr lang
\end{proof}

\begin{lem}[Extremallemma]
	Für jede Alternative $y$ gilt: Wenn jeder Wähler die Alternative $y$ als bester (\glqq Top \grqq) oder als letzter (\glqq Flop \grqq) gewählt hat, dann muss die kollektive Auswahlfkt die Alternative $y$ ebenfalls als \glqq Top \grqq oder \glqq Flop \grqq setzen. 
\end{lem}

\begin{rem}[Alternative Formulierung des Satzes von Arrow]
	\textbf{demokratische Grundregeln:}
	\begin{itemize}
		\item Auswahlfkt K muss auf ganz $\mathcal{P}_A$ formuliert sein
		\item Ergebnis von K muss in $\mathcal{P}_A$ liegen
		\item Die Pareto Bedingung muss erfüllt sein
		\item die Unabhängigkeit von irrelevanten Alternativen muss erfüllt sein
		\item es gibt keinen Diktator
	\end{itemize}
	
	Eine Alternative Formulierung von dem Satz von Arrow ist nun: Es kann keine kollektive Auswahlfunktion geben, die alle demoktratischen Grundregeln erfüllt. 
\end{rem}

\section{Nicht-Manipulierbkeit von Wahlen}
Das Problem beim Satz von Arrow ist es, dass das Einfügen von irrelevcanten Alternativen zwischen bevorzugten Kandidaten und Hauptkonkurrenten kann das Ergebnis beeinflussen. 

\begin{defi}[soziale Entscheidungsfkt]
	$e: \mathcal{P}_A^n \rightarrow A$
\end{defi}

\begin{defi}[Strategische Manipulation/ Anreizkompatibel]
	Eine soziale Entscheidungsfunktion $e: \mathcal{P}_A^n \rightarrow A$ kann durch Wähler strategisch manipuliert werden, falls es Präferenzen $\rho_1, \dots, \rho_i, \dots, \rho_n, \rho_i' \in \mathcal{P}_A$ gibt, sodass $\rho_i \neq \rho_i'$ 
	\begin{align*}
		a \rho_i b \text{ gilt für }  a=e(\rho_1, \dots, \rho_i, \dots, \rho_n), b=e(\rho_1, \dots, \rho_i', \dots, \rho_n)
	\end{align*} 
	Die Fkt $e$ heißt Anzeizkompatibel, falls $e$ nicht strategisch manipulierbar ist. 
\end{defi}

\begin{rem}[zu definition Strategische Manipulation]
	e-manipulierbar, falls ein Wähler $i$, der $b$ vor $a$ präferiert, $b$ erzwingen kann, wenn er statt seiner Präferenz $\rho_i$ eine Präferenz $\rho_i' \neq \rho_i$ angibt.  
\end{rem}

\begin{defi}[Monotonie]
	Eine soziale Entscheidungsfkt $e$ heißt monoton, falls aus 
	$ a=e(\rho_1, \dots, \rho_i, \dots, \rho_n), b=e(\rho_1, \dots, \rho_i', \dots, \rho_n)$ mit $a \neq b$ folgt, dass $a \rho_i b$ und $a \rho_i' b$ gilt. 
\end{defi}

\begin{thm}
	Eine soziale Entscheidungsfunktion ist genau dann monoton, wenn sie Anreizkompatibel ist. 
\end{thm}

\begin{defi}[Diktator eine Entscheidungsfkt]
	Wähler $i$ heißt Diktator einer sozialen Entscheidungsfkt, falls für alle $\rho_1, \dots \rho_n \in \mathcal{P}_A$ gilt, dass $ a=e(\rho_1, \dots, \rho_n)$, wobei a der eindeutig bestimmte Kandidat mit $a \rho_i b$ für alle $b \neq a$ ist. Die Fkt $e$ heißt diktatorisch, falls er einen Diktator besitzt. 
\end{defi}

\begin{defi}[Top Menge]
	Sei $S \subset A$ und $\rho \in \mathcal{P}_A$. Wir führen eine Präferenzrelation $\rho^S$ ein. für 
	\begin{align*}
		a,b \in S: a \rho^s b \Leftrightarrow s \rho b \\
		a,b \notin S: a \rho^s b \Leftrightarrow s \rho b \\
		a \in S, b \notin S; a \rho^S b
	\end{align*}
\end{defi}
Man kann zeigen, dass $\rho^s$ dadurch eindeutig bestimmt ist. 

\begin{lem}[Top Präferenz]
	Sei $e$ eine Anreizkompatible und surjektive Entscheidungsfkt. Dann gilt: 
	\begin{align*}
		\forall \rho_i, \dots \rho_n \in \mathcal{P}_A \text{ und } S \subset A, S \neq \emptyset \text{ gilt: } e(\rho_q^S, \dots, \rho_n^S) \in S
	\end{align*}
\end{lem}

\begin{defi}[Erweiterung der Entscheidungsfkt]
	Die Funktion $E: \mathcal{P}_A^n \rightarrow \mathcal{P}_A$ ist die Erweiterung der sozialen Entscheidungsfkt e und ist definiert durch $E(\rho_1, \dots, \rho_n)$ wobei 
	\begin{align*}
		a \rho b : \Leftrightarrow e(\rho_1^{\{a,b\}} , \dots, \rho_n^{\{a,b\}})=a \forall a,b \in A
	\end{align*}
\end{defi}

\begin{lem}
	Falls e eine Anreizkompatible und surjektive soziale Entscheidungsfkt ist, dann ist ihre Erweiterung eine kollektive Auswahlfkt. 
\end{lem}

\begin{lem}[Erweiterungslemma]
	Falls e eine anreizkompatible, surjektive und nicht diktatorische Entscheidungsfkt ist, so ist ihre Erweiterung E eine kollektive Auswahlfkt, die Einstimmigkeit, Unabhängigkeit von irrelevanten Ereignissen und nicht diktator.
\end{lem}

\begin{thm}[Satz von Gibbord-Salterthwaite]
	Falls $e$ eine surjektive, anreizkompatible Entscheidungsfkt ist, sodass drei oder mehr Alternativen wählbar sind, dann ist e diktatorisch. 
\end{thm}







% --- Hauptteil der Arbeit II
% ========================
% Hauptteil der Arbeit II
% ========================

\chapter{Informationssuche im Netz}





% --- Zusammenfassung, Fazit und Ausblick
%% ====================================
% Zusammenfassung, Fazit und Ausblick
% ====================================

\chapter{Fazit und Ausblick}
\label{chap:fazit_ausblick}



\end{document}
